\section*{Chapter \ref{ch:g10:soil-formation} --- Soil, a Slow-Forming Resource}

\begin{solution}{exo:g10:soil-formation:1}
\omterm{def:g10:soil-formation:chemical-weathering}{Chemical weathering} is the transformation of the minerals of a rock by
reactions with water and the gases it carries (carbon dioxide, oxygen): some
minerals dissolve, others become new minerals stable at the surface.
\omterm{def:g10:soil-formation:hydrolysis}{Hydrolysis} is the \omterm{def:g10:soil-formation:chemical-weathering}{chemical weathering} of a silicate by water and its hydrogen
ions: the metal ions go into solution and the rest of the mineral is rebuilt
into a \omterm{def:g10:soil-formation:clay}{clay mineral}.
\end{solution}

\begin{solution}{exo:g10:soil-formation:2}
\omterm{def:g10:soil-formation:weathering-profile}{Saprolite} is rock decomposed in place: soft and earthy, it keeps the grains,
bands and structure of its \omterm{def:g10:soil-formation:weathering-profile}{parent rock} and passes downwards, without a break,
into the fresh rock. A soil formed in river sand rests on material brought
from elsewhere: its grains are rounded and sorted, and it shows none of the
structure of the rock beneath it.
\end{solution}

\begin{solution}{exo:g10:soil-formation:3}
The feldspars become clay (kaolinite), releasing potassium, sodium or
calcium and silica into the water; the black mica becomes clay and iron
oxides; the quartz is left unchanged, as sand.
\end{solution}

\begin{solution}{exo:g10:soil-formation:4}
The iron of the dark minerals, \ce{Fe^2+}, is oxidised by the oxygen carried
by the water into insoluble iron(III) oxides: yellow-brown goethite and red
hematite, which coat the grains and colour the whole material.
\end{solution}

\begin{solution}{exo:g10:soil-formation:5}
\qty{0.017}{mm} a year and \qty{1.537}{mm} a year: $1.537/0.017 \approx 90$
times larger.
\end{solution}

\begin{solution}{exo:g10:soil-formation:6}
Calcium: 1 and 1. Carbon: $1 + 1 = 2$ and 2. Oxygen: $3 + 2 + 1 = 6$ and
$2 \times 3 = 6$. Hydrogen: 2 and 2. Charge: 0, and $+2 - 2 = 0$.
\end{solution}

\begin{solution}{exo:g10:soil-formation:7}
The dissolution of calcite produces only dissolved ions, \ce{Ca^2+} and
\ce{HCO3-}: everything is carried away and a pure limestone leaves no solid
residue. The \omterm{def:g10:soil-formation:hydrolysis}{hydrolysis} of feldspar produces a solid, kaolinite, which stays;
and the quartz of the granite hardly reacts at all. Granite therefore leaves
clay and sand.
\end{solution}

\begin{solution}{exo:g10:soil-formation:8}
The clay flakes and the humus carry negative charges that hold the positive
potassium ions \ce{K+} on their surfaces, so that rainwater passing through
does not carry them away. A root releases hydrogen ions, which take the place
of the potassium on the flakes; the potassium set free in the soil water is
taken up by the root.
\end{solution}

\begin{solution}{exo:g10:soil-formation:9}
\omterm{def:g10:soil-formation:chemical-weathering}{Chemical weathering} needs water, abundant in wet climates and scarce in dry
ones (and frozen in cold ones); and its reactions run faster when it is
warm. Only under heat and abundant water, over a long time, do the reactions
go far enough to remove almost everything soluble and leave iron and
aluminium oxides.
\end{solution}

\begin{solution}{exo:g10:soil-formation:10}
$10/0.017 \approx 590$ years for \qty{1}{cm}; $300/0.017 \approx 17\,600$
years for \qty{30}{cm}; $1000/0.017 \approx 59\,000$ years for \qty{1}{m}.
\end{solution}

\begin{solution}{exo:g10:soil-formation:11}
Volume: $10\,000 \times 0.0015 = \qty{15}{m^3}$; mass: $15 \times 1200 =
\qty{18000}{kg}$, \qty{18}{tonnes} per hectare per year.
\end{solution}

\begin{solution}{exo:g10:soil-formation:12}
$46\,000/0.017 \approx 2.7$ million years. Such a thick \omterm{def:g10:soil-formation:weathering-profile}{saprolite} needs a
landscape that has been stable, without its weathered mantle being stripped
away, for millions of years.
\end{solution}

\begin{solution}{exo:g10:soil-formation:13}
Goethite: \ce{4 Fe^2+ + O2 + 6 H2O -> 4 FeO(OH) + 8 H+}: iron 4 and 4,
oxygen $2 + 6 = 8$ and $4 \times 2 = 8$, hydrogen 12 and $4 + 8 = 12$, charge
$+8$ and $+8$. Hematite: \ce{4 Fe^2+ + O2 + 4 H2O -> 2 Fe2O3 + 8 H+}: iron 4
and 4, oxygen $2 + 4 = 6$ and 6, hydrogen 8 and 8, charge $+8$ and $+8$.
\end{solution}

\begin{solution}{exo:g10:soil-formation:14}
Net loss $0.082 - 0.017 = \qty{0.065}{mm}$ a year; $300/0.065 \approx 4600$
years. Under conventional ploughing, $300/1.52 \approx 200$ years: more than
twenty times shorter.
\end{solution}

\begin{solution}{exo:g10:soil-formation:15}
The tonne contains \qty{300}{kg} of orthoclase. Each kilogram gives
\qty{464}{g} of kaolinite: $300 \times 0.464 \approx \qty{139}{kg}$. Potassium:
$300 \times 0.1405 \approx \qty{42}{kg}$.
\end{solution}

\begin{solution}{pb:g10:soil-formation:1}
\textbf{1.} $4 \times 0.3048 \approx \qty{1.22}{m}$.

\textbf{2.} $0.5 \times 30.48 \approx \qty{15}{cm}$.

\textbf{3.} Humus, the dark organic matter made by the decomposers from dead
plants.

\textbf{4.} Iron oxides, goethite and hematite, formed by the oxidation of
the iron(II) of the minerals of the parent material.

\textbf{5.} No. Below, the profile would pass into weathered rock, the
\omterm{def:g10:soil-formation:weathering-profile}{saprolite} (or other weathered material), and then into the fresh \omterm{def:g10:soil-formation:weathering-profile}{parent
rock}.

\textbf{6.} \ce{2 KAlSi3O8 + 2 H+ + 9 H2O -> Al2Si2O5(OH)4 + 2 K+ + 4 H4SiO4}.

\textbf{7.} $1000/278.33 \approx \qty{3.59}{mol}$.

\textbf{8.} $3.59/2 \approx \qty{1.80}{mol}$ of kaolinite, $1.80 \times 258.16
\approx \qty{464}{g}$.

\textbf{9.} The potassium ions and the dissolved silica: they are carried by
the water to the rivers, and by the rivers to the sea.

\textbf{10.} It hardly reacts: the quartz grains are left as sand.

\textbf{11.} $1219/0.017 \approx 72\,000$ years.

\textbf{12.} $1.537 - 0.017 = \qty{1.52}{mm}$ a year.

\textbf{13.} $1219/1.52 \approx 800$ years.

\textbf{14.} $72\,000/800 = 90$: the profile is lost ninety times faster than
it formed.

\textbf{15.} No: $0.013 < 0.017$, erosion is a little slower than production
and the soil holds its own.

\textbf{16.} $T = S/(E - P)$.

\textbf{17.} $300/1.52 \approx 200$ years.

\textbf{18.} $300/(0.082 - 0.017) = 300/0.065 \approx 4600$ years.

\textbf{19.} A soil lost in a few centuries takes tens of thousands of years
to form again: on the scale of human lives and societies, what is lost is
not replaced.

\textbf{20.} $1000/1.52 \approx 660$ years: \textbf{a metre of soil is lost
in about 660 years, a few centuries, under conventional ploughing}.
\end{solution}
